%; whizzy paragraph -pdf xpdf -latex ./whizzypdfptex.sh
%; whizzy-paragraph "^\\\\begin{frame}\\|\\\\emtext"
% latex beamer presentation.
% platex, latex-beamer でコンパイルすることを想定。 

%     Tokyo Debian Meeting resources
%     Copyright (C) 2012 Junichi Uekawa

%     This program is free software; you can redistribute it and/or modify
%     it under the terms of the GNU General Public License as published by
%     the Free Software Foundation; either version 2 of the License, or
%     (at your option) any later version.

%     This program is distributed in the hope that it will be useful,
%     but WITHOUT ANY WARRANTY; without even the implied warreanty of
%     MERCHANTABILITY or FITNESS FOR A PARTICULAR PURPOSE.  See the
%     GNU General Public License for more details.

%     You should have received a copy of the GNU General Public License
%     along with this program; if not, write to the Free Software
%     Foundation, Inc., 51 Franklin St, Fifth Floor, Boston, MA  02110-1301 USA

\documentclass[cjk,dvipdfmx,12pt]{beamer}
\usetheme{Tokyo}
\usepackage{monthlypresentation}

%  preview (shell-command (concat "evince " (replace-regexp-in-string "tex$" "pdf"(buffer-file-name)) "&")) 
%  presentation (shell-command (concat "xpdf -fullscreen " (replace-regexp-in-string "tex$" "pdf"(buffer-file-name)) "&"))
%  presentation (shell-command (concat "evince " (replace-regexp-in-string "tex$" "pdf"(buffer-file-name)) "&"))

%http://www.naney.org/diki/dk/hyperref.html
%日本語EUC系環境の時
\AtBeginDvi{\special{pdf:tounicode EUC-UCS2}}
%シフトJIS系環境の時
%\AtBeginDvi{\special{pdf:tounicode 90ms-RKSJ-UCS2}}

\newenvironment{commandlinesmall}%
{\VerbatimEnvironment
  \begin{Sbox}\begin{minipage}{1.0\hsize}\begin{fontsize}{8}{8} \begin{BVerbatim}}%
{\end{BVerbatim}\end{fontsize}\end{minipage}\end{Sbox}
  \setlength{\fboxsep}{8pt}
% start on a new paragraph

\vspace{6pt}% skip before
\fcolorbox{dancerdarkblue}{dancerlightblue}{\TheSbox}

\vspace{6pt}% skip after
}
%end of commandlinesmall

\title{東京エリアDebian勉強会}
\subtitle{第122回 2015年1月度}
\author{野島貴英}
\date{2015年1月17日}
\logo{\includegraphics[width=8cm]{image200607/openlogo-light.eps}}

\begin{document}

\begin{frame}
\titlepage{}
\end{frame}

\begin{frame}{設営準備にご協力ください。}
会場設営よろしくおねがいします。
\end{frame}

\begin{frame}{Agenda}
 \begin{minipage}[t]{0.45\hsize}
  \begin{itemize}
   \item 注意事項
	 \begin{itemize}
	  \item 写真はセミナールーム内のみ可です。
          \item 出入りは自由でないので、もし外出したい方は、野島まで一声くださいませ。
	 \end{itemize}
   \item 事前課題発表
  \end{itemize}
 \end{minipage} 
 \begin{minipage}[t]{0.45\hsize}
  \begin{itemize}
   \item 最近あったDebian関連のイベント報告
	 \begin{itemize}
	  \item 第121回 東京エリアDebian勉強会
	 \end{itemize}
   \item Debian Trivia Quiz
   \item Emacs関連パッケージのDebianパッケージ作成
   \item 今後のイベント
   \item 今日の宴会場所
  \end{itemize}
 \end{minipage}
\end{frame}

\section{事前課題}
\emtext{事前課題}
{\footnotesize
\begin{prework}{ $BLnEg(B }
  \begin{enumerate}
  \item Q.hack time$B$K2?$r$7$^$9$+!)(B\\
    A. RC$B%G%P%C%0!"(BDebian unstable$B>e$NI>2A!"%P%0%l%]=q$-$J$I!#(BRC$B$O(B\url{https://bugs.debian.org/release-critical/}$B;2>H!#(B
  \item ($B%*%W%7%g%s(B)Q.$B2?$K$D$$$FJ9$-$?$$!?;22C<T$HOC$r$7$?$$$G$9$+!)(B\\
    A. $B<+M3%=%U%H%&%'%"$H<+M3$J4D6-$H<+M3$J5;=Q$K$D$$$F$NG.$$;W$$!*<+M3%=%U%H%&%'%"$H<+M3$J4D6-$H5;=Q$G<R2q$X2?$r%"%T!<%k$7$?J}$,$"$J$?$K$H$C$F(BHappy$B$+!"$=$7$F!"M_K>$r!"$$$?$:$i?4$rK~$?$;$k$N$+!#(B
  \end{enumerate}
\end{prework}

\begin{prework}{ alohaug }
  \begin{enumerate}
  \item Q.hack time$B$K2?$r$7$^$9$+!)(B\\
    A. debian livecd for gnuk key generation
  \item ($B%*%W%7%g%s(B)Q.$B$I$3$G:#2s$NJY6/2q$N3+:E$rCN$j$^$7$?$+!)(B\\
    A. $B$=$NB>(B
  \item ($B%*%W%7%g%s(B)Q.$B2?$K$D$$$FJ9$-$?$$!?;22C<T$HOC$r$7$?$$$G$9$+!)(B\\
    A. gnuk$B$dHkL)80J]4I%*!<%W%s%O!<%I$N:r:#(B
  \end{enumerate}
\end{prework}

\begin{prework}{ keiichirou\_miyamoto }
  \begin{enumerate}
  \item Q.hack time$B$K2?$r$7$^$9$+!)(B\\
    A. Linux$B$NJY6/Ey(B
  \end{enumerate}
\end{prework}

\begin{prework}{ henrich }
  \begin{enumerate}
  \item Q.hack time$B$K2?$r$7$^$9$+!)(B\\
    A. $BB?J,869F=q$$$F$^$9!D!#(B
  \item ($B%*%W%7%g%s(B)Q.$B$I$3$G:#2s$NJY6/2q$N3+:E$rCN$j$^$7$?$+!)(B\\
    A. $BM'C#$dCN$j9g$$$+$iD>@\(B
  \item ($B%*%W%7%g%s(B)Q.$B2?$K$D$$$FJ9$-$?$$!?;22C<T$HOC$r$7$?$$$G$9$+!)(B\\
    A. $B3'$,(BDebian$B$K4X$7$FCN$j$?$$$3$H$rGD0.$7$?$$$G$9(B
  \end{enumerate}
\end{prework}

\begin{prework}{ koedoyoshida }
  \begin{enumerate}
  \item Q.hack time$B$K2?$r$7$^$9$+!)(B\\
    A. $B$^$C$?$j%a%s%F%J%s%9$J$I!#(B
  \item ($B%*%W%7%g%s(B)Q.$B$I$3$G:#2s$NJY6/2q$N3+:E$rCN$j$^$7$?$+!)(B\\
    A. Debian JP$B$N%a!<%j%s%0%j%9%H(B
  \end{enumerate}
\end{prework}

\begin{prework}{ dictoss }
  \begin{enumerate}
  \item Q.hack time$B$K2?$r$7$^$9$+!)(B\\
    A. Debian GNU/kFreeBSD$B$r%M%?$K$7$?%W%l%<%s%F!<%7%g%s=`Hw(B
  \item ($B%*%W%7%g%s(B)Q.$B$I$3$G:#2s$NJY6/2q$N3+:E$rCN$j$^$7$?$+!)(B\\
    A. Debian JP$B$N%a!<%j%s%0%j%9%H(B
  \end{enumerate}
\end{prework}

\begin{prework}{ kenhys }
  \begin{enumerate}
  \item Q.hack time$B$K2?$r$7$^$9$+!)(B\\
    A. RC$B%P%0$X$NBP1~(B
  \item ($B%*%W%7%g%s(B)Q.$B$I$3$G:#2s$NJY6/2q$N3+:E$rCN$j$^$7$?$+!)(B\\
    A. Debian JP$B$N%a!<%j%s%0%j%9%H(B
  \end{enumerate}
\end{prework}


}

\section{イベント報告}
\emtext{イベント報告}

\begin{frame}{第121回東京エリアDebian勉強会}

\begin{itemize}
\item 場所はスクウェア・エニックスさんのセミナルームをお借りしての開催でした。
\item 参加者は9名でした。
\item セミナ内容は 
\begin{itemize}
\item @kenhysさんにより「DebianとFedoraでパッケージをリリースするまでの話」
\item 野島さんにより、「Debianから見たLinux Mint」
\end{itemize}
でした。
\item 残りの時間でhack timeを行い、成果発表をしました。
\item 宴会の代わりに、「大戸屋ごはん処 新宿イーストサイドスクエア」で夕食会をやりました。
\end{itemize} 
  
\end{frame}

\begin{frame}{第121回東京エリアDebian勉強会(つづき)}

　セミナですが、@kenhysさんにより、groongaについての概要とDebian/Fedoraへ公式パッケージをリリースするまでが語られました。現在のDebian/Fedoraの公式パッケージ化に関して、実際のところが語られました。実際に公式パッケージにするまでの諸々について語られる機会は、他ではあまり見ない件なので、非常に貴重な発表だったと思います。また、野島さんにより、Linux Mintについて、Debianとの比較について語られました。
  
\end{frame}


\section{Debian Trivia Quiz}
\emtext{Debian Trivia Quiz}
\begin{frame}{Debian Trivia Quiz}

  Debian の常識、もちろん知ってますよね?
知らないなんて恥ずかしくて、知らないとは言えないあんなことやこんなこと、
みんなで確認してみましょう。

今回の出題範囲は\url{debian-devel-announce@lists.debian.org},
\url{debian-news@lists.debian.org} に投稿された
内容などからです。

\end{frame}

\subsection{問題}

%; whizzy-master ../debianmeetingresume201311.tex
% $B0J>e$N@_Dj$r$7$F$$$k$?$a!"$3$N%U%!%$%k$G(B M-x whizzytex $B$9$k$H!"(Bwhizzytex$B$,MxMQ$G$-$^$9!#(B
%

\santaku
{12/29$B$K!"(BDebian$B$N%5!<%S%9$HO"7H$G$-$k(Bchromium$B8~$1$N3HD%5!G=$,%j%j!<%9$5$l$^$7$?!#2?$HO"7H$G$-$k$N$G$7$g$&$+!)(B}
{tracker.debian.org}
{sources.debian.net}
{www.debian.or.jp}
{B}
{zack$B$5$s$N(BDebconf14$B$N%W%l%<%s$r8+$F%$%s%9%Q%$%"$5$l$?$H$N$3$H$G!"(BRaphael Geissert$B$5$s$,:n$C$?(Bchromium$BMQ(B(chrome$B$G$bF0$/!K$N%*%s%i%$%s%=!<%9%3!<%I%(%G%#%?$r<B8=$9$k3HD%5!G=$G$9!#(BDebian$B%Q%C%1!<%8$N%=!<%9%3!<%I$rD>@\%*%s%i%$%s$G1\Mw$G$-$k(Bsources.debian.net$B>e$N%3!<%I$r$3$N3HD%5!G=$r;H$C$F%7!<%`%l%9$KJT=82DG=$K$J$j$^$9!#(Bzack$B$5$s$N<($7$??<9o$JLdBj!JBh(B117$B2sEl5~%(%j%"(BDebian$BJY6/2q;qNA;2>H!K$r2r7h$9$k0Y$N2h4|E*$J%"%$%G%"$N#1$D$G$9!*%^%B$9$P$i$7$$!*<+M3%=%U%H%&%'%"%i%V!"(BDebian$B%i%V$J?M$O!"AaB.!"(Bchromium$B$X(BNautilus$B$+$i$3$N3HD%5!G=$rEj$29~$s$G;n$9$Y$7!*(B}

\santaku
{2014/11/11$B$K(BRapha\"el Hertzog$B$5$s$K$h$jDs0F$5$l$?(BDEP14$B$NFbMF$O0J2<$N$I$l!)(B}
{git$B%Q%C%1!<%8%s%0%j%]%8%H%j$N;H$$J}$K$D$$$F(B}
{control$B%U%!%$%k$K(BDebian-Homepage$B%U%#!<%k%I$rF~$l$k7o(B}
{debian/upstream/metadata$B$rMQ0U$9$k7o(B}
{A}
{$B%Q%C%1!<%8%s%0%j%]%8%H%j$N%V%i%s%A$H%?%0$K%Y%s%@!<>pJs!JNc!'(Bdebian/*,ubuntu/*)$B$d$i%P!<%8%g%s>pJs$NF~$lJ}$J$I$r4^$`!"=t!9$N%k!<%k$NDs0F!#>\$7$/$O!"(Bhttp://dep.debian.net/deps/dep14/$B;2>H!#$J$*!"(BDebian-Homepage$B%U%#!<%k%I$N7o$O(BDEP13$B!#(Bdebian/upstream/metadata$B$N7o$O(BDEP12$B!#(B}

\santaku
{2015/1/8 Marvel$B<R$h$j(BDebian ARM$BMQ%Q%C%1!<%8%S%k%I%^%7%s$N4sIU$,$"$j$^$7$?!#$5$F2?Bf$NDs6!$,$"$C$?$G$7$g$&!)(B}
{7$BBf(B}
{8$BBf(B}
{10$BBf(B}
{B}
{Marvel$B<R$,Ds6!$7$?%5!<%P$O(B MV78460 SoC$B%Y!<%9$N%\!<%I(B8$B$D$rDs6!$H$N$3$H!#$J$*!"(BMV78460$B$O!"(BARM v7$B%3%"#44p$rEk:\$9$k(BMarvel$B<R3+H/$N(BSoC$B!#;29M!'(Bhttp://www.marvell.com/embedded-processors/armada-xp/}

\santaku
{2015/1/10$B$K(Bwheezy$B$,%"%C%W%G!<%H$5$l$^$7$?!#$5$F%P!<%8%g%s$O<!$N$I$l!#(B}
{7.6}
{7.7}
{7.8}
{C}
{wheezy$B$r$*;H$$$NJ}!9$OAaB.%"%C%W%0%l!<%I$7$^$7$g$&!#(B}

\santaku
{2015/1/1$B$K$"$kD9$5L$K~$N(BDebian keyring$B$N>C5n$,40N;$7$?$=$&$G$9!#$"$kD9$5$H$O(B?}
{2048}
{4096}
{8192}
{A}
{$B%"%J%&%s%9$H!">C$5$l$?3:Ev$N%-!<0lMw$O(Bdebian-devel-announce$B$N(B2015/1/1$B$N5-;v$K5-:\$,$"$j$^$9!#(Bhttps://lists.debian.org/debian-devel-announce/2015/01/msg00000.html}

\santaku
{2014/12/6$B$K(Bcdn.debian.net$B$N%l%3!<%I$N<($9@h$r!"$H$"$k(BFQDN$B$N(BCNAME$B$H$7$?$$;]$NLd$$9g$o$;$,(Bdebian-devel ML$B$K$"$j$^$7$?!#$I$3$K8~$1$k!)(B}
{ftp.debian.or.jp}
{ftp.jp.debian.org}
{http.debian.net}
{C}
{cdn.debian.net$B$rDs6!$7$F$$$k(BAraki$B$5$s$+$i$NLd$$9g$o$;!#(Bsqueeze$B0J>e$GEk:\$5$l$F$$$k(Bapt ver 0.7.21$B$G$O(Bhttp$B$N(Bredirect$B$KBP1~$7$F$$$k$N$G!"MxMQ<T$K:G$b6a$$%Q%C%1!<%8%@%&%s%m!<%I@h$r65$($F$/$l$k5!9=$r(Bhttp.debian.net$B$K0lK\2=$7$?$$$H$N$3$H!#FC$KH?BP$9$k?M$O5o$J$$LOMM$G$9!#(B}

\santaku
{2015/1/9$B$K(Bdebian-devel ML$B$K$F!"0lIt%Q%C%1!<%8$N(Bdescription$B$,D9$9$.$k7o$K$D$$$F$N5DO@$,$"$j$^$7$?!#:G$bD9$$(Bdescription$B$N9T$r;}$D%Q%C%1!<%8$O<!$N$I$l!)(B}
{firmware-linux-nonfree}
{texlive-latex-extra}
{irssi-scripts}
{B}
{2015/1/9$B8=:_!"(Bunstable/main$B$G$O!":G$bD9$$9T?t$N=g$G$$$/$H!"(B1$B0L(B texlive-latex-extra($B<B$K(B1935$B9T(B!!)$B!"(B2$B0L(B texlive-fonts-extra(437$B9T(B)$B!"(B3$B0L(B irssi-scripts(350$B9T(B)$B!#(Bunstable/non-free$B$G$O!"(B1$B0L(B firmware-linux-nonfree(251$B9T(B)$B$G$9!#(B}




\section{Emacs関連パッケージのDebianパッケージ作成}
\emtext{Emacs関連パッケージのDebianパッケージ作成}

\section{今後のイベント}
\emtext{今後のイベント}
\begin{frame}{今後のイベント}
\begin{itemize}
 \item 関西エリアDebian勉強会
 \item 東京エリアDebian勉強会　
\end{itemize}
\end{frame}

\section{今日の宴会場所}
\emtext{今日の宴会場所}
\begin{frame}{今日の宴会場所}
未定
\end{frame}

\end{document}

;;; Local Variables: ***
;;; outline-regexp: "\\([ 	]*\\\\\\(documentstyle\\|documentclass\\|emtext\\|section\\|begin{frame}\\)\\*?[ 	]*[[{]\\|[]+\\)" ***
;;; End: ***
